\begin{lemma}[Step 1, Subcase 2.2 with a positive divider]
Let $G,e_0,[k],T,\varphi$ be a finite loopless multigraph with exact edge
identities, distinguished uncolored edge, color set, positive-rung
certified ETT, and generating coloring.  Assume maximum property,
properness, exact one-uncoloredness, and the complete six-part lower-rung
package.  Let $H$ be a good hierarchy over a hierarchy base, with
$d<q_0$ dividers.  Assume the same-rung fewer-than-$q_0$ elementarity
hypothesis for arbitrary certified inputs and generating colorings, the
large bound $|\overline\varphi(T_n)|\geq2n+11$, the separate terminal
certificate
\[
 T\text{ is }
 \left(\bigcup_{\eta\in D^H_{n,d+1}}\Gamma^d_\eta\right)^-
 \text{-closed under }\varphi,
\]
and elementarity of $V(T)$.  Put $A=H.\mathit{starred}(d)$.

Let $\alpha\ne\beta$, with $\alpha$ missing and singleton-closed on $T$
and $\beta$ not missing on $T$.  Let $P_1,P_2,P_3$ be pairwise
support-disjoint exact $(T,\varphi,\{\alpha,\beta\})$-exit paths.  Their
outside ends $a_1,a_2,a_3$ and tree ends $b_1,b_2,b_3$ are pairwise
distinct, with $b_1\prec_T b_2\prec_T b_3$.  Retain the exact
$\beta$-colored boundary-edge identity $f_j$ of each path incident with
$b_j$.  Suppose $b_3\in V(A)$.

Assume that a selected color $\theta\ne\alpha,\beta$ is missing and
singleton-closed on $A$ and singleton-closed on $T$.  Assume also the
explicit no-smaller-presentation consequence of the global minimum: for
every proper, exactly one-uncolored coloring $\sigma$ stable from
$\varphi$ on the final ladder tree, arbitrary replacement input,
corresponding same-rung maximum ETT $T'$, good hierarchy with fewer than
$d$ dividers, its own large estimate, and its separate terminal
minus-closed certificate, there do not exist two support-disjoint exact
exit legs with distinct outside ends for a missing singleton-closed first
color.  Assume additionally that $d>0$.  Then the displayed Subcase 2.2
data are contradictory.
\end{lemma}
\begin{proof}
Write $n$ for the rung number, $T_n$ for the last ladder tree, $D_i$ for
the exact cumulative stored-color set at slot $i$, and $\varphi_i$ for
the exact coloring at that slot.  Thus $\varphi=\varphi_n$ by the
generating-coloring field of \noderef{ExtendedTashkinovTree}.  We first
make the three exit legs used in every branch completely explicit.

The starred tree $A$ is an ordered initial segment of $T$.  Unfold
\noderef{TashkinovHierarchy}.  If $d>0$, then
\noderef{HierarchyStarredStage} identifies $A$ with the $d$-th hierarchy
stage, and repeated use of the nestedness field of the hierarchy gives
the required initial-segment map to its terminal stage $T$.  If $d=0$,
the same conclusion follows from the core-inside-first field and then the
nestedness and terminal-stage fields of the hierarchy.  All these are
entrywise statements by \noderef{TreeSequenceInitialSegment}.
Unfolding \noderef{TreeVertexOrder} gives positions $r_1<r_2$ displaying
$b_1,b_2$ and positions $s_2<s_3$ displaying $b_2,b_3$.  The vertex map
of a \noderef{TreeSequence} is injective, so its two occurrences of
$b_2$ are identical: $r_2=s_2$.  Next unfold
\noderef{TreeSequenceVertexSet} in $b_3\in V(A)$ and choose a position
$t_3$ of $A$ displaying $b_3$.  The vertex-entry clause of the
initial-segment relation makes the canonical image of $t_3$ a position
of $T$ displaying $b_3$.  Vertex injectivity identifies that image with
the occurrence $s_3$ supplied by the second order relation.  Hence
$s_3\leq A.\mathit{length}$.  The inequalities
$r_1<r_2=s_2<s_3$ put the numerical positions $r_1,r_2$ inside
$A$; applying the entrywise equality at those two positions shows that
they display $b_1,b_2$ in $A$.  Together with the assumed membership of
$b_3$, this proves $b_1,b_2,b_3\in V(A)$.

Unfold each given \noderef{TreeExitPath}.  The edge-node immediately
preceding the terminal vertex-node $b_j$ represents an exact boundary
edge of $T$: its other adjacent vertex-node on the path is outside $T$,
because $b_j$ is the only vertex-node of the path in $V(T)$.  By
\noderef{TwoColorIncidenceGraph} that edge has color $\alpha$ or
$\beta$.  It cannot have color $\alpha$, since
\noderef{TreeColorSetClosed} says that the exact $\alpha$-colored
boundary set is empty.  Hence this last boundary edge is
$\beta$-colored.  The separately retained identity $f_j$ has the same
two properties by hypothesis; unfolding \noderef{BoundaryColorEdges}
and \noderef{BoundaryEdges} shows in particular that its end other than
$b_j$ lies outside $T$.  Since $b_j\in V(A)\subseteq V(T)$, the same
identity $f_j$ is also a $\beta$-colored boundary edge of $A$.

We next justify which $\alpha,\theta$ components may be exchanged.  If an
$\alpha$-edge node occurs on $P_j$, it is not incident with $T$:
otherwise one of its incident vertex-nodes on this alternating path
would be in $T$, hence would be $b_j$, and the edge would be an
$\alpha$-colored boundary edge, contrary to the preceding closedness.
More generally, its full \noderef{TwoColorComponent} for
$\alpha,\theta$ contains no vertex of $T$.  For if it did, take a
shortest incidence-graph walk from the external edge to a vertex-node in
$T$ and its first vertex-node in $T$.  The preceding exact edge identity
crosses the boundary of $T$ and, by
\noderef{TwoColorIncidenceGraph}, has color $\alpha$ or $\theta$.
This contradicts respectively the assumed singleton closedness of
$\alpha$ or of $\theta$ on $T$.  The identical first-crossing argument
shows that every $\alpha,\theta$ component containing an outside vertex
is wholly outside $T$.

Exchange $\alpha$ and $\theta$ on every full $\alpha,\theta$ component
outside $T$, and call the result $\nu$.  Finiteness makes this a finite
choice of roots, and the preceding paragraph verifies the outside-support
clause of \noderef{OutsideKempeRecoloring}.  By
\noderef{OutsideKempeRecoloringStable}, $\nu$ is proper, leaves exactly
$e_0$ uncolored, fixes every exact edge incident with $T$, and is
$(T,C,\varphi_n)$-stable for every chosen control set $C$.  We verify its
restriction to $T_n$.  The above-penultimate field of
\noderef{ExtendedTashkinovTree} and
\noderef{TreeSequenceInitialSegment} give the canonical entrywise
inclusion $V(T_n)\subseteq V(T)$ after unfolding
\noderef{TreeSequenceVertexSet}.  If an exact edge $f$ is incident with
$T_n$, unfold \noderef{TreeIncidentEdge} and retain its witnessing
endpoint; the vertex inclusion makes the same endpoint witness that $f$
is incident with $T$.  Likewise, unfolding \noderef{TreeMissingColors}
as the finite union of the vertex missing-color sets shows
\[
 \overline{\varphi_n}(T_n)\subseteq
 \overline{\varphi_n}(T):
\]
a witness vertex in $V(T_n)$ is sent by the same inclusion to $V(T)$,
with its \noderef{MissingColors} membership unchanged.  Now unfold
\noderef{StableColoring}.  For its edge clause on $T_n$, transport the
incident-edge witness as just described; a controlling color in
$\overline{\varphi_n}(T_n)\cup C$ lies in
$\overline{\varphi_n}(T)\cup C$ by the displayed inclusion, so the edge
clause on $T$ fixes $f$.  For its vertex clause, transport the vertex to
$V(T)$ and apply the vertex clause there.  Thus both clauses hold and
$\nu$ is $(T_n,C,\varphi_n)$-stable for every $C$.

For $j=1,2,3$, form a walk $Q_j$ with the same ordered vertex- and
edge-nodes as $P_j$.  Each old $\beta$-edge remains $\beta$ because
$\beta\ne\alpha,\theta$.  Every old $\alpha$-edge lies in one of the
selected outside components just proved and becomes $\theta$.  Thus
every successive pair of nodes of $P_j$ is adjacent in
$I_G(\nu;\theta,\beta)$, so $Q_j$ is a walk there; it is the same simple
node list and has the same positive length and support as $P_j$.
Its outside end still misses one of its two colors.  Explicitly, if
$\beta$ was missing at $a_j$, then the definition
\noderef{MissingColors}, together with \noderef{PresentColors}, says
that no incident edge had color $\beta$; none acquires that color in an
$\alpha,\theta$ exchange, so $\beta$ remains missing.  If instead
$\alpha$ was missing, then no incident edge had color $\alpha$.  The
whole $\alpha,\theta$ component through the outside vertex $a_j$ was
selected, so the $\theta$-edges incident with $a_j$ become $\alpha$ and
the new $\theta$-edges are exactly its old $\alpha$-edges; hence
$\theta$ is missing there under $\nu$.  We already proved that all three
tree ends lie in $A$, while $a_j\notin V(T)$ implies
$a_j\notin V(A)$.  Finally, the only vertex-node of $Q_j$ in $A$ is
$b_j$, because the corresponding claim for $P_j$ excludes every other
vertex even from the larger tree $T$.  This checks every field of
\[
       \operatorname{TreeExitPath}(A,\nu,\theta,\beta,a_j,b_j,Q_j).
       \tag{1}
\]
The supports remain pairwise disjoint, the three outside ends remain
pairwise distinct, and no claim about separation of the three
transformed full components has been made.

Apply \noderef{HierarchyLEMStable} to the just constructed outside
recoloring.  It supplies a corresponding ETT $\mathcal T^\nu$ displayed
on $T$, a hierarchy base $B^\nu$, a hierarchy $H^\nu$, and a good
hierarchy on it.  The corresponding certificate retains the rung, all
ordered ladder trees, exact connecting-edge sets, stored sets, and kinds;
the same-hierarchy conclusion retains every ordered stage, divider,
stored-color set, and Gamma pair.  The result also supplies maximum
property and transports the given terminal minus-closed certificate.
The stability established above identifies the missing-color sets on
$T_n$ under $\nu$ and $\varphi_n$, so the large bound is unchanged.

By the additional hypothesis, $d>0$.  In $H^\nu$, the tree $A$ is again the exact
$d$-th stage.  Apply \noderef{GoodHierarchyPrefixRepackage} with $r=d$.
It returns an ETT $\mathcal E_A$ displayed on the exact ordered tree
$A$, having generating coloring $\nu$, the same rung and literally the
same dependent certified series as $\mathcal T^\nu$, together with
maximum property, a base $B_A$, a hierarchy $K$ with $d-1$ dividers,
and all five \noderef{GoodHierarchy} fields.

We now construct, rather than assume, the corresponding certificate for
$\mathcal E_A$ relative to the original $\mathcal T$.  Take its
replacement input, replacement series, and displayed ETT to be those of
$\mathcal E_A$.  The prefix theorem identifies that series with the
series of $\mathcal T^\nu$.  Compose this equality, slot by slot, with
the certificate already carried by
\noderef{CorrespondingExtendedTashkinovTree} for $\mathcal T^\nu$.
This proves equality with the original series of every retained ordered
ladder tree, every exact stored set, every exact connecting-edge set,
and every extension kind.  It also transports unchanged the properness,
one-uncoloredness, and exact stable-coloring fields at all replacement
slots.  The prefix theorem gives the same rung and generating coloring
$\nu$.  These are precisely all fields of a corresponding ETT; call the
resulting certificate $\mathcal C_A$.

The remaining hypotheses of the no-smaller-presentation clause also hold
with their exact indices.  First, its stability is the
$(T_n,D_n,\varphi_n)$-stability of $\nu$ already obtained directly from
\noderef{OutsideKempeRecoloringStable}; this refers to the original final
ladder tree, not merely to the new displayed tree $A$.  Second,
$\mathcal E_A$ has the same series and rung, so its last ladder tree is
the original $T_n$, and equality of its missing sets under $\nu$ and
$\varphi_n$ proves the unchanged large bound.  Third, unfold
\noderef{HierarchyTerminalStoredColors}.  Since $K$ has $d-1$ dividers,
its terminal stored set is $D_n\setminus\overline\nu(A)$.  Unfolding
\noderef{HierarchyStoredColorsAt} at the old index $d$ gives exactly the
same set, because the old $d$-th starred stage is $A$.  The Gamma at the
last index of $K$ is the old Gamma at $d-1$ by the prefix theorem.
Consequently its separately returned closure
\[
 A\text{ is }
 \left(\bigcup_{\eta\in D^{H^\nu}_{n,d}}\Gamma^{d-1}_\eta\right)^-
 \text{-closed under }\nu
\]
is literally the terminal minus-closed expression demanded for $K$.
Finally $d-1<d$, $\theta$ is still missing and singleton-closed on $A$
because the first recoloring fixes every edge incident with $A$ and every
missing set there, and the paths $Q_1,Q_2$ from (1) have disjoint
supports and different outside ends.  Applying the stated
no-smaller-presentation hypothesis to
$\nu,A,\mathcal C_A,B_A,K$ now gives a contradiction.

\end{proof}
